\documentclass{article}
\usepackage[margin=1in]{geometry}
\usepackage{amsmath, amssymb, amsthm, enumitem, multicol}

\everymath{\displaystyle}
\newtheorem*{theorem*}{Theorem}
\newtheorem*{definition*}{Definition}
\newcommand{\vct}[1]{\langle #1 \rangle}
\newcommand{\vu}{\mathbf{u}}
\newcommand{\vv}{\mathbf{v}}
\newcommand{\vw}{\mathbf{w}}
\newcommand{\vx}{\mathbf{x}}
\newcommand{\vn}{\mathbf{n}}

\newcommand{\colvecTwo}[2]{
  \begin{bmatrix}
    #1 \\
    #2
  \end{bmatrix}
}
\newcommand{\colvecThree}[3]{
  \begin{bmatrix}
    #1 \\
    #2 \\
    #3
  \end{bmatrix}
}

\newcommand{\casesTwo}[2]{
  \begin{cases}
    #1 \\
    #2
  \end{cases}
}
\newcommand{\casesThree}[3]{
  \begin{cases}
    #1 \\
    #2 \\
    #3
  \end{cases}
}

\title{Linear Algebra: A Modern Introduction}
\author{David Poole}
\date{}

\begin{document}
\maketitle
\setcounter{section}{1}
\section{Systems of Linear Equations}

\subsection{Introduction to Systems of Linear Equations}
\begin{definition*}
  A linear equation in the $n$ variables $x_1, x_2, \ldots, x_n$ is an equation that can be writen
  in the form
  \[a_1x_1+a_2x_2+\cdots+a_nx_n=b\]
  where the coefficients $a_1, a_2, \ldots, a_n$ and the constant term $b$ are constants.
\end{definition*}
\subsubsection*{Exercises}
\begin{enumerate}
    \begin{multicols}{3}
    \item Yes
    \item No
    \item No
    \item No
    \item No
    \item Yes
    \end{multicols}
  \item $2x+4y=7$
  \item $x+y=1, x \ne y$
  \item $x+y=4, x\ne0, y\ne0$
  \item $x-100y=0, x>0, y>0$
    \setcounter{enumi}{38}
  \item
    \begin{enumerate}
      \item $2x+y=3$
      \item $x=\frac{3}{2}-\frac{1}{2}s, y=s$
    \end{enumerate}
    \setcounter{enumi}{40}
  \item
    \begin{align*}
      2t+3s & = 0 \\
      3t+4s & = 1
    \end{align*}
    $x=\frac{1}{3},y=2\frac{1}{2}$
\end{enumerate}

\subsection{Direct Methods for Solving Linear Systems}
\begin{definition*}
  A matrix is in \textbf{row echelon form} if it satisfies the following properties:
  \begin{enumerate}
    \item Any rows consisting entirely of zeroes are at the bottom.
    \item In each nonzero row, the first nonzero entry (called the \textbf{leading entry}) is in a
      column to the left of any leading entries below it.
  \end{enumerate}
\end{definition*}
\begin{definition*}
  The following \textbf{elementary row operations} can be performed on a matrix:
  \begin{enumerate}
    \item Interchange two rows.
    \item Multiply a row by a nonzero constant.
    \item Add a multiple of a row to another row.
  \end{enumerate}
\end{definition*}
\begin{definition*}
  Matrices $A$ and $B$ are \textbf{row equivalent} if there is a sequence of elementary row
  operations that converts $A$ into $B$.
\end{definition*}
\begin{theorem*}
  Matrices $A$ and $B$ are row equivalent if and only if they can be reduced to the same row echolon
  form.
\end{theorem*}
\begin{proof}
  If $A$ and $B$ are row equivalent, the further row operations will reduce $B$ (which $A$ can be
  converted to) to some (same) row echelon form.

  Conversely, if $A$ and $B$ have the same row echelon form $R$, then, via elementary row
  operations, we can convert $A$ into $R$ and $B$ into $R$. Reversing the latter sequence of
  operations, we can convert $R$ into $B$, and therefore the sequence $A\to R\to B$ achieves the
  desired effect.
\end{proof}
\begin{definition*}
  The \textbf{rank} of a matrix is the number of nonzero rows in its row echelon form.
\end{definition*}
\begin{theorem*}[The Rank Theorem]
  Let $A$ be the coefficient matrix of a system of linear equations with $n$ variables. If the system
  is consistent, then
  \[\text{numer of free variables}=n-\operatorname{rank}(A).\]
\end{theorem*}
\end{document}
